\section{Introduction}\label{sec:introduction}
Sentiment analysis identifies opinions expressed in text. In this project, an
English movie review is mapped to a binary label: negative (0) or positive (1).
The application is relevant to summarizing audience feedback, but even this
two-class formulation requires more than counting positive and negative words.
A reviewer may praise the acting and criticize the plot, reverse a judgment with
negation, or use apparently positive language sarcastically.

We investigate how text representation affects this task under a shared data
split. TF-IDF encodes lexical evidence, including short phrases; Word2Vec
compresses a review into a dense average of word vectors; DistilBERT produces
context-sensitive representations and is adapted using LoRA. The classical
branch evaluates Logistic Regression and a linear Support Vector Machine with
each of its two representations, producing four classical pipelines. Adding
DistilBERT + LoRA gives five complete pipelines in total.

The experimental questions are:
\begin{enumerate}
  \item How do TF-IDF and averaged Word2Vec compare with the same classifiers?
  \item Does LoRA-adapted DistilBERT improve the recorded test performance, and
        how do its input length and adaptation settings affect validation F1?
  \item Which linguistic phenomena remain difficult, and how can an interactive
        demonstration expose differences between the models?
\end{enumerate}

The implementation covers collection and quality auditing, fixed split creation,
three motivated preprocessing techniques, representation learning, classifier
training, quantitative evaluation and qualitative error analysis. All models use
the same review IDs and labels. Classical models consume filtered, lemmatized
tokens, while the transformer consumes minimally cleaned natural text. Thus the
comparison concerns complete pipelines, including pretrained knowledge and tuning
effort; it is not an isolated ablation of representation or preprocessing.

The report follows the eight sections requested in the course assignment.
Section~\ref{sec:related} summarizes relevant work;
Section~\ref{sec:data} describes the data; Sections~\ref{sec:methods}
and~\ref{sec:experiments} explain the implementation and experiments;
Sections~\ref{sec:results} and~\ref{sec:errors} analyze performance and errors;
Section~\ref{sec:conclusion} concludes. Appendices provide task allocation,
commit evidence and reproduction instructions.
